\pageoriginale

\clearpage
\SGAExposeBookmark{XIX}{El teorema de Noether}
\begin{center}
{\large SGA 7 --- EXPOSÉ XIX}\par
\vspace{2em}
{\Large\underline{El teorema de Noether}}\par
\vspace{1.5em}
{\large por P. Deligne}
\end{center}

\section*{Sumario}

\noindent 1. Enunciado de los teoremas. \hfill 1

\noindent 2. Un complemento a XVIII. \hfill 4

\noindent 3. Demostración de 1.3. \hfill 5

\noindent 4. Un segundo complemento a XVIII. \hfill 8

\noindent 5. Los casos excepcionales del teorema de Noether. \hfill 11

\section*{1. Enunciado de los teoremas}

La presente exposición actualiza la demostración de Lefschetz
[3, p.~108] del teorema de Noether (1.2) enunciado a continuación.

\subsubsection*{1.1.}

Sea $k$ un cuerpo algebraicamente cerrado de característica $p$. Sean
$d\geq1$ y
\[
\underline a=(a_i)_{1\leq i\leq d}
\]
una sucesión de enteros mayores o iguales que $1$. Una intersección
completa $V$ de dimensión $n$ y multigrado $\underline a$ en el espacio
proyectivo $\mathbb P^{n+d}$ sobre $k$ se dirá \underline{genérica} si,
tras una extensión de $k$, es proyectivamente isomorfa a una
intersección completa definida por ecuaciones cuyos coeficientes son
algebraicamente independientes sobre el cuerpo primo. En lo que sigue,
salvo indicación expresa en contrario, suponemos $d\geq1$ y $a_i\geq2$.

\subsubsection*{Teorema 1.2.}

Sea $V\subset\mathbb P^{2+d}$ una superficie intersección completa
genérica de multigrado $\underline a$. Supongamos que
\[
\underline a\ne(2),(3),(2,2).
\]
Entonces toda curva reducida de $V$ es la intersección esquemática de
$V$ con una hipersuperficie de $\mathbb P^{2+d}$.

Mostremos que el teorema resulta de su corolario

\pageoriginale

\begin{equation}\tag{1.2.1}
\operatorname{Pic}(\mathbb P^{2+d})=\mathbb Z
\xrightarrow{\ \sim\ }\operatorname{Pic}(V).
\end{equation}
Si $C$ es una curva de $V$, (1.2.1) implica
$\mathcal O(C)\simeq\mathcal O(n)$ para cierto $n$. Si $H$ es una
sección hiperplana, entonces
\[
n=\bigl(\textstyle\prod_i a_i\bigr)^{-1}(C,H)>0.
\]
Queda probar que el sistema lineal inducido sobre $V$ por $|nH|$ es
completo. Escribamos $V$ como intersección de hipersuperficies $H_i$ de
grado $a_i$, y pongamos $V_i=H_1\cap\cdots\cap H_i$. Una inducción sobre
$i$ muestra que el sistema lineal inducido por $|nH|$ sobre $V_i$ es
completo: en la sucesión exacta
\[
H^0(V_i,\mathcal O(n))\longrightarrow
H^0(V_{i+1},\mathcal O(n))\longrightarrow
H^1(V_i,\mathcal O(n-a_{i+1})),
\]
el último grupo es nulo (FAC).

\subsubsection*{Teorema 1.3.}

Sea $V\subset\mathbb P^{2n+d}$ una intersección completa genérica de
dimensión par no nula $2n$ y de multigrado $\underline a$. Supongamos
que se verifica una de las condiciones siguientes:
\begin{enumerate}[label=(\alph*)]
\item $k$ no es de característica $2$, y $(2n;\underline a)$ es
  distinto de $(2n;2)$, $(2n;2,2)$ y $(2;3)$;
\item el número de Hodge $h^{2n,0}$ no es nulo, es decir,
  $\sum_i(a_i-1)>2n$.
\end{enumerate}
Entonces toda clase de cohomología algebraica de
$H^{2n}(V,\mathbb Q_\ell(n))$ es un múltiplo racional de la clase
$\eta^n$ de una sección lineal.

Deducimos (1.2.1) de 1.3. Los casos excepcionales de 1.2 son exactamente
aquellos para los que
$\sum_i(\sisi{d_i}{a_i}-1)\leq2$; por tanto, 1.3(b) se aplica a la
superficie $V$ de 1.2. Sabemos que $\operatorname{Pic}^{\tau}(V)=0$
(XI 1.8), luego $\operatorname{Pic}(V)$ se inyecta en
$H^2(V,\mathbb Q_\ell(1))$. Por 1.3, todo elemento de
$\operatorname{Pic}(V)$ es un múltiplo racional de la clase $\eta$ de
$\mathcal O(1)$, y (XI 1.8) da la conclusión.

Es claro que la validez de 1.2 y 1.3 no depende de la intersección
completa genérica elegida, sino únicamente de $p$, $n$ y
$\underline a$.

\pageoriginale

\subsubsection*{Teorema 1.4.}

Sea $X\subset\mathbb P$ una variedad proyectiva lisa, pura de dimensión
$2n+1$ sobre $k$. Sea $\ell\subset\check{\mathbb P}$ un pincel de
Lefschetz de secciones hiperplanas de $X$. Si $k$ tiene característica
$2$, supongamos además que $\ell$ es un pincel de Lefschetz genérico.

Supongamos que $X$ satisface la condición (LV) (XVIII 5.2.2), y sea $V$
la sección hiperplana genérica del pincel $\ell$. Se verifica una de las
dos condiciones siguientes:
\begin{enumerate}[label=(\alph*)]
\item ninguna clase de cohomología no nula de la parte evanescente
  $\operatorname{Ev}(V,\mathbb Q_\ell(n))$ de
  $H^{2n}(V,\mathbb Q_\ell(n))$ es algebraica;
\item las clases de cohomología algebraica engendran
  $\operatorname{Ev}(V,\mathbb Q_\ell(n))$.
\end{enumerate}

Precisemos el sentido de ``genérico'' en 1.4.
\begin{enumerate}[label=(\alph*)]
\item Sea $G$ la grassmanniana de rectas del espacio proyectivo
  $\check{\mathbb P}$ dual de $\mathbb P$. Por ``punto geométrico'' de
  $G$ entendemos un morfismo del espectro de un cuerpo algebraicamente
  cerrado a $G$. Todo punto geométrico $\bar s\to G$ define un pincel
  $\ell$ de secciones hiperplanas de la variedad obtenida de $X$
  extendiendo escalares de $k$ a $k(\bar s)$. Decimos que $\ell$ es
  genérico si $\bar s$ está sobre el punto genérico de $G$.
\item Sea $\ell$ un pincel de secciones hiperplanas (una recta de
  $\check{\mathbb P}$). Todo punto geométrico $\bar s\to\ell$ define
  una sección hiperplana de la variedad obtenida de $X$ extendiendo
  escalares de $k$ a $k(\bar s)$. Decimos que esta sección es genérica
  si $\bar s$ está sobre el punto genérico de $\ell$.
\end{enumerate}

Sea $\bar s$ como arriba, con imagen el punto genérico $s$ de $\ell$, y
sea $S\subset\ell$ el conjunto finito de puntos que dan secciones
hiperplanas singulares. El grupo de monodromía
$\pi_1(\ell-S,\bar s)$, cociente de
$\operatorname{Gal}(k(\bar s)/k(s))$, actúa sobre
$\operatorname{Ev}(V,\mathbb Q_\ell(n))$. Si $k$ no es de
característica $2$, esta acción es irreducible (XVIII 6.7); en la
sección 2 mostraremos que la afirmación sigue siendo válida sin
hipótesis sobre la característica para un pincel genérico. El
subespacio de $\operatorname{Ev}(V,\mathbb Q_\ell(n))$ engendrado por
las clases de cohomología algebraica es estable bajo monodromía, es
decir, bajo Galois. Es, pues, cero o todo el espacio, lo que prueba 1.4.

\pageoriginale

\subsubsection*{Corolario 1.5.}

Sea $V\subset\mathbb P^{2n+d}$ una intersección completa genérica de
dimensión par $2n$ y multigrado $\underline a$. Entonces:
\begin{enumerate}[label=(\alph*)]
\item toda clase de cohomología algebraica de
  $H^{2n}(V,\mathbb Q_\ell(n))$ es un múltiplo racional de $\eta^n$; o
\item $H^{2n}(V,\mathbb Q_\ell(n))$ está engendrado por clases de
  cohomología algebraica.
\end{enumerate}
Si se verifica (b), entonces, para toda intersección completa de
dimensión $2n$ y multigrado $\underline a$ sobre un cuerpo
algebraicamente cerrado de característica $p$, el grupo
$H^{2n}(V,\mathbb Q_\ell(n))$ está engendrado por clases algebraicas.

Apliquemos 1.4 con $X$ una intersección completa genérica de dimensión
$2n+1$ y multigrado $(a_1,\ldots,a_{d-1})$, y con $\ell$ un pincel
genérico de secciones de $X$ por hipersuperficies de grado $a_d$
(cf. XVII 2.4). También podríamos tomar $X$ de multigrado
$(a_1,\ldots,a_d)$ y $\ell$ un pincel genérico de secciones
hiperplanas. La afirmación final resulta de que una clase de cohomología
algebraica sigue siendo algebraica por especialización.

\section*{2. Un complemento a XVIII}

\subsubsection*{Proposición 2.1.}

Sea $X\subset\mathbb P$ una variedad proyectiva lisa, pura de dimensión
\sisi{par $2n$}{impar $2n+1$} sobre un cuerpo algebraicamente cerrado
$k$, y sea $\ell\subset\check{\mathbb P}$ un pincel de Lefschetz
genérico de secciones hiperplanas de $X$. Sea $s$ un punto geométrico
de $\ell$, sea $H_s$ la sección hiperplana correspondiente, sea
$S\subset\ell$ el conjunto de valores excepcionales del pincel, y sea
$\pi_1(\ell-S,s)$ el grupo de monodromía. Supongamos que $X$ satisface
la condición (LV). Entonces la representación de $\pi_1(\ell-S,s)$
sobre la cohomología evanescente
\[
\operatorname{Ev}(H_s,\mathbb Q_\ell)
\subset H^{2n}(H_s,\mathbb Q_\ell)
\]
es absolutamente irreducible.

Recordemos que $\operatorname{Ev}(H_s,\mathbb Q_\ell)$ es el ortogonal
de la imagen de $H^{2n}(X,\mathbb Q_\ell)$ en
$H^{2n}(H_s,\mathbb Q_\ell)$, imagen que coincide con el subgrupo de
invariantes bajo el grupo de monodromía (XVIII 5.6.10).

\pageoriginale

\subsubsection*{Demostración.}

La recta proyectiva $\mathbb P^1$ es simplemente conexa. Por tanto,
$\pi_1(\ell-S,s)$ está engendrado topológicamente por los subgrupos de
inercia en los puntos de $S$ y sus conjugados. Su imagen en
\[
\operatorname{GL}\bigl(
\sisi{H^n(H_s,\mathbb Q_\ell)}{H^{2n}(H_s,\mathbb Q_\ell)}
\bigr)
\]
está engendrada topológicamente por las transformaciones
\[
x\longmapsto x-(-1)^n(x\cdot\delta)\delta,
\]
donde $\delta$ es conjugado de un ciclo evanescente. En consecuencia,
la imagen de $H^{2n}(X,\mathbb Q_\ell)$ en
$H^{2n}(H_s,\mathbb Q_\ell)$ es el ortogonal del subespacio engendrado
por los ciclos evanescentes y sus conjugados; así,
$\operatorname{Ev}(H_s,\mathbb Q_\ell)$ está engendrado por dichos
ciclos y sus conjugados.

Para un pincel de Lefschetz \underline{genérico}, los argumentos de
XVIII, §6, muestran que los ciclos evanescentes asociados a los
distintos puntos de $S$ son conjugados entre sí.

Sea $W\subset\operatorname{Ev}(H_s,\mathbb Q_\ell)$ un subespacio
invariante no nulo. Existen un conjugado $\delta$ de un ciclo
evanescente y un elemento $v\in W$ tales que $(v,\delta)\ne0$; aquí es
donde se utiliza (LV). Puesto que
\[
v-(-1)^{\sisi{m}{n}}(v\cdot\delta)\delta\in W,
\]
se tiene $\delta\in W$. El subespacio $W$, que contiene un ciclo
evanescente y todos sus conjugados, coincide por tanto con
$\operatorname{Ev}(H_s,\mathbb Q_\ell)$. El argumento sigue siendo
válido tras extender el cuerpo de coeficientes $\mathbb Q_\ell$, lo que
prueba la irreducibilidad absoluta.

\subsubsection*{Observación 2.2.}

Si se omite la condición (LV), los argumentos precedentes muestran aún
que todo subespacio invariante de $H^{2n}(H_s,\mathbb Q_\ell)$ contiene
$\operatorname{Ev}(H_s,\mathbb Q_\ell)$ o está contenido en la imagen
de $H^{2n}(X,\mathbb Q_\ell)$. La misma afirmación vale para $X$ de
dimensión par y para todo pincel de Lefschetz.

\section*{3. Demostración de 1.3}

\subsubsection*{3.1.}

Bajo las hipótesis del teorema queda probar que no se verifica 1.5(b).
El caso $h^{0,2n}\ne0$ será tratado por Katz en XXI 4.1; la condición
$h^{0,2n}\ne0$ equivale a
\[
\sum_i\sisi{d_i}{a_i}-(2n+r+1)\geq0,
\]
y este número es el entero $m$ tal que
$\Omega_V^{2n}\simeq\mathcal O(m)$. Los resultados de la presente exposición
no volverán a utilizarse en el resto del seminario, de modo que no hay
círculo vicioso.

\pageoriginale

\subsubsection*{3.2.}

Cuando $k=\mathbb C$, se demuestra 1.3 observando que, si se verifica
1.5(b), la cohomología de $V$ es puramente de tipo $(n,n)$, puesto que
la cohomología algebraica es de tipo $(n,n)$. Esto solo ocurre en los
casos excepcionales excluidos en 1.3 (XI 2.9). Por el principio de
Lefschetz, 1.3 es, pues, verdadero en característica cero.

\subsubsection*{Proposición 3.3.}

Supongamos las hipótesis de 1.4 y que se verifica (b). Sea
$\operatorname{Ev}(V,\mathbb Q)$ el subespacio vectorial sobre
$\mathbb Q$ de $\operatorname{Ev}(V,\mathbb Q_\ell(n))$ engendrado por
las clases de cohomología algebraica.
\begin{enumerate}[label=(\alph*)]
\item Se tiene un isomorfismo
  \[
  \operatorname{Ev}(V,\mathbb Q)\otimes\mathbb Q_\ell
  \xrightarrow{\ \sim\ }
  \operatorname{Ev}(V,\mathbb Q_\ell(n)).
  \]
  Sobre $\operatorname{Ev}(V,\mathbb Q)$, la forma de intersección es
  racional y definida, de signo $(-1)^n$.
\item Los ciclos evanescentes y sus conjugados pertenecen a
  $\operatorname{Ev}(V,\mathbb Q)$; forman allí un sistema de raíces
  irreducible de tipo $A$, $D$ o $E$.
\item El grupo de monodromía, imagen de $\pi_1(\ell-S,\bar s)$ en
  $\operatorname{GL}(\operatorname{Ev}(V,\mathbb Q))$, es finito; es
  el \sisi{grupo de Weil}{grupo de Weyl} del sistema de raíces de (b).
\end{enumerate}

Salvo el carácter definido de la forma de intersección, (a) resulta de
que la forma de intersección es no degenerada sobre
$\operatorname{Ev}(V,\mathbb Q_\ell(n))$ por (LV), y racional sobre
$\operatorname{Ev}(V,\mathbb Q)$.

La condición 1.4(b) equivale a que existan sobre $V$ ciclos algebraicos,
módulo equivalencia numérica, que no sean invariantes por monodromía;
por tanto, esta condición no depende de $\ell$. Sea $x$ tal ciclo
algebraico, y sea $\delta$ el ciclo evanescente asociado a
$s_o\in S$ y a un camino $c$ desde $\bar s$, que define $V$, hasta
$s_o$. Para un $c$ adecuado, $(x,\delta)\ne0$. El transformado de
Galois
\[
\sigma(x)=x-(-1)^n(x,\delta)\delta
\]
es aún algebraico; por tanto, $(x,\delta)\delta$ es algebraico, imagen
de un ciclo $y$ independiente de $\ell$. El número racional
\[
\frac{(-1)^n}{2}(y,y)=(x,\delta)^2
\]
es, en consecuencia, independiente de $\ell$. Es un cuadrado en todo
$\mathbb Q_\ell$, luego es un cuadrado en $\mathbb Q$; así, salvo
signo, $(x,\delta)$ pertenece a $\mathbb Q$ y es independiente de
$\ell$. Esto muestra no solo que
$\delta\in\operatorname{Ev}(V,\mathbb Q)$, sino también que el ciclo
algebraico módulo equivalencia numérica
\[
\pm\delta=\frac{\sigma(x)-x}{(x,\delta)}
\]
es independiente de $\ell$.

\pageoriginale

Un ciclo algebraico sobre una variedad $V$ está siempre definido sobre
una extensión finita de todo cuerpo de definición de $V$. Se sigue que
$\pi_1(\ell-S,s)$ actúa sobre $\operatorname{Ev}(V,\mathbb Q)$ a través
de un grupo finito. Este grupo está engendrado por las reflexiones
asociadas a los ciclos evanescentes $\delta$ y sus conjugados
(XVIII 6.6.1 y §2). Por tanto, los $\delta$ y sus conjugados forman un
sistema de raíces cuyo \sisi{grupo de Weil}{grupo de Weyl} es el grupo
de monodromía. Además, todas las raíces de este sistema son conjugadas
bajo el \sisi{grupo de Weil}{grupo de Weyl}; es de tipo $A$, $D$ o $E$.

Salvo un factor escalar, la forma de intersección es la única forma
invariante bajo el \sisi{grupo de Weil}{grupo de Weyl}. Es, pues,
definida. Su signo se determina calculando
\[
(\delta,\delta)=2(-1)^n.
\]

\subsubsection*{Proposición 3.4.}

Si $k=\mathbb C$, bajo las hipótesis de 1.4 son equivalentes las
condiciones siguientes:
\begin{enumerate}[label=(\roman*)]
\item el grupo de monodromía es finito;
\item la cohomología evanescente $\operatorname{Ev}(V)$ es puramente
  de tipo $(n,n)$.
\end{enumerate}

Sea $\widetilde X\to\ell$ la familia cuyas fibras son las secciones
hiperplanas del pincel $\ell$, y sea $V_s$ la fibra en
$s\in\ell(\mathbb C)$. Los $\operatorname{Ev}(V_s)$ forman una
variación algebraica de estructuras de Hodge sobre $\ell$. Por el
teorema de la parte fija [2, 4.1.2], aplicado a un recubrimiento finito
de $\ell$ (cf. [2, 4.1.3.3]), (i) implica que la estructura de Hodge
sobre $\operatorname{Ev}(V_s)$ es localmente constante. El grupo de
monodromía actúa, por tanto, mediante automorfismos de la estructura de
Hodge. Para una transformación de Picard--Lefschetz $T$,
$\operatorname{Im}(T-1)$ es una subestructura de Hodge de rango uno;
es, pues, de tipo $(n,n)$ y está engendrada por un ciclo evanescente
$\delta$. Como los ciclos evanescentes son de tipo $(n,n)$, todo
$\operatorname{Ev}(V_s)$ es de ese tipo.

Recíprocamente, si $\operatorname{Ev}(V)$ es puramente de tipo $(n,n)$,
la forma de intersección sobre $\operatorname{Ev}(V)$ es definida y el
grupo de monodromía es finito (cf. [2, 4.1.3.3]).

\pageoriginale

\subsubsection*{Observación 3.5.}

Si se verifican las condiciones equivalentes de 3.3, el argumento de
3.3 muestra que los ciclos evanescentes y sus conjugados forman un
sistema de raíces de tipo $A$, $D$ o $E$, cuyo
\sisi{grupo de Weil}{grupo de Weyl} es el grupo de monodromía.

\subsubsection*{3.6.}

Demostremos 1.3 cuando $k$ tiene característica finita distinta de
$2$. Sea $S=\operatorname{Spec}(V)$ el espectro de
un anillo de valoración discreta de característica mixta con cuerpo
residual $k$, y sea
$X\subset\mathbb P_S^{2n+d}$ una intersección completa de dimensión
relativa $2n+1$ y multigrado $(a_1,\ldots,a_{d-1})$ sobre $S$. Sea
$\ell$ un pincel de Lefschetz de secciones de $X$ por
hipersuperficies de grado $a_d$; en la fibra especial, $\ell_s$ es un
pincel de Lefschetz para $X_s$. Sea
\[
f:\widetilde X\longrightarrow\ell\simeq\mathbb P^1_S
\]
la familia cuyas fibras son las secciones del pincel. El lugar
excepcional $T\subset\ell$ es étale sobre $S$ (XVII 6; aquí se usa
$p\ne2$). El haz $R^{2n}f_*\mathbb Q_\ell$ es localmente constante
sobre $\ell-T$; por tanto, está moderadamente ramificado y su
monodromía es la misma en característica $p$ que en característica
cero. Si se verifica 1.5(b), se deduce de 3.3(c) y 3.4 que el grupo de
cohomología $H^{2n}(V)$ de las intersecciones completas complejas de
dimensión $2n$ y multigrado $\underline a$ es puramente de tipo
$(n,n)$, cosa que solo ocurre en los casos excluidos por 1.3
(XI 2.9).

\section*{4. Un segundo complemento a XVIII}

En esta sección trabajamos sobre $\mathbb C$ y con cohomología clásica
entera. Retomamos, con métodos más próximos a los de Lefschetz, algunos
resultados de XVIII. Allí solo se demostraron después de tensorizar por
$\mathbb Q$ (o, más precisamente, por $\mathbb Q_\ell$, lo que viene a
ser lo mismo).

\subsubsection*{4.1.}

Sea $X\subset\mathbb P$ una variedad proyectiva compleja lisa, pura de
dimensión $n+1$, y sea $\ell$ un pincel de Lefschetz de secciones
hiperplanas de $X$, con eje $A$. Denotamos asimismo por $\ell$ el
espacio de parámetros del pincel, isomorfo a $\mathbb P^1$. Sea
$S\subset\ell$ el conjunto de valores excepcionales. Para $s\in\ell$,
denotemos por $H_s$ la sección hiperplana correspondiente de $X$, de
dimensión $n$.

\pageoriginale

Sea $T(A)$ un entorno tubular cerrado de $A\cap X$ en $X$, cuyo borde
es transversal a todos los $H_s$.

Elijamos dos puntos $0,\infty\in\ell-S$, y sea
\[
\alpha:\sisi{\ell-S}{\ell}\longrightarrow\mathbb P^1(\mathbb C)
\]
un difeomorfismo que preserva la orientación, envía $0$ en $0$,
$\infty$ en $\infty$ y $S$ en la circunferencia de radio uno centrada
en $0$. El compuesto $\alpha f$ de $\alpha$ con la proyección sobre
$\ell$ hace de $X-T(A)^\circ$ una fibración singular en variedades con
borde sobre $\mathbb P^1(\mathbb C)$, con una singularidad sobre cada
punto de $\alpha(S)$. Denotemos por $c_s$ el segmento rectilíneo que une
$0$ con $\alpha(s)$.

Omitiendo $\alpha$ de la notación, se obtiene la figura
\begin{center}
\begin{tikzpicture}[x=1.2cm,y=1.2cm]
  \coordinate (O) at (0,0);
  \coordinate (L) at (-1.8,1.8);
  \coordinate (S2) at (0,2.8);
  \coordinate (S1) at (2.2,2.0);
  \draw (O)--(L);
  \draw (O)--(S2);
  \draw (O)--(S1);
  \node[right] at (O) {$0$};
  \node[above right] at (S2) {$s_2$};
  \node[above right] at (S1) {$s_1$};
  \node[right] at (0,1.65) {$c_2$};
  \node[below right] at (1.25,1.14) {$c_1$};
  \fill (-0.73,0.76) circle (0.035);
  \fill (-0.55,0.91) circle (0.035);
  \fill (-0.37,1.03) circle (0.035);
\end{tikzpicture}
\end{center}

Cerca de $s\in S$ y del punto singular $x_s$ de $H_s\subset X$, unas
coordenadas locales adecuadas nos sitúan en la situación estudiada en
XIV 3.2. Si $\alpha$ se elige holomorfa cerca de $s$, incluso pueden
elegirse coordenadas locales sobre $X$ y $\ell$ en las que $f$ se
escribe $\sum z_i^2$ y $\alpha^{-1}(c_s)$ viene dado por ``los $z_i$
son reales y positivos''. Sobre la intersección de un entorno de $s$
con $\alpha^{-1}(c_s)-\{s\}$, consideremos el fibrado de esferas reales
$\Delta'_s\subset X$ definido por ``los $z_i$ son reales''. Cuando
$t\to s$, la fibra $\Delta'_{st}$ de $\Delta'_s$ en $t$ tiende a
$x_s$.

Sea $\Delta_s$ la unión de $\{x_s\}$ y de un fibrado de esferas sobre
$\alpha^{-1}(c_s)-\{s\}$, contenido en $X-T(A)$ y que prolonga
$\Delta'_s$. Entonces $\Delta_s$ es un disco encajado en $X$, cuyo
borde $H_0\cap\Delta_s$ es una esfera que representa el ciclo
evanescente $\delta_s$ en $s$, transportado hasta $H_0$ a lo largo del
camino $\alpha^{-1}(c_s)$.

Apliquemos la teoría de Morse a la función $|\alpha f|^2$ para estudiar
cómo cambia el tipo de homotopía de
\[
\{x\in X\mid x\in T(A)\ \text{o bien}\ |\alpha f(x)|^2\leq r\}
\]

\pageoriginale

cuando $r$ crece. Pongamos
\[
X^-=\{x\in X\mid x\in T(A)\ \text{o bien}\ \alpha f(x)\ne\infty\}.
\]
Se obtiene que las inclusiones
\[
H_0\cup\bigcup_{s\in S}\Delta_s
\lhook\joinrel\longrightarrow
T(A)\cup H_0\cup\bigcup_{s\in S}\Delta_s
\lhook\joinrel\longrightarrow X^-
\]
son equivalencias de homotopía. La teoría de Morse muestra asimismo que
$X-X^-$ se retrae por deformación sobre un subcomplejo de dimensión
real a lo sumo $n$. En conjunto, esto da el resultado siguiente.

\subsubsection*{Teorema 4.2.}

Existe un subcomplejo compacto $\Sigma$ de dimensión real a lo sumo
$n$ en $H_\infty^{\mathrm{af}}=H_\infty-A$ tal que la inclusión
\[
H_0\cup\bigcup_{s\in S}\Delta_s\longrightarrow X-\Sigma
\]
es una equivalencia de homotopía.

\subsubsection*{4.3.}

La supresión de $\Sigma$ no afecta a los grupos de cohomología,
homología u homotopía de $X$ en grados $i\leq n$. Como
$H_0\cup\bigcup_{s\in S}\Delta_s$ se obtiene de $H_0$ adjuntando discos
de dimensión $n+1$ a lo largo de esferas que representan los ciclos
evanescentes, la homología da
\[
\mathbb Z^S\longrightarrow
H_n\bigl(\sisi{V_0}{H_0}\bigr)
\longrightarrow H_n(X)\longrightarrow0.
\]
La sucesión exacta análoga en cohomología, junto con la sucesión que se
deduce de ella por dualidad de Poincaré, forma el diagrama conmutativo
\begin{equation}\tag{4.3.1}
\begin{tikzcd}[row sep=large,column sep=large]
{} & \mathbb Z^S & {} & {} \\
\mathbb Z^S \arrow[r,"{\sum n_s\delta_s}"] &
H^n\bigl(\sisi{V_0}{H_0}\bigr)
  \arrow[u,"{(\ ,\delta_s)}"'] \arrow[r] &
H^{n+2}(X) \arrow[r] & 0 \\
{} & H^n(X) \arrow[u] \arrow[ur,"{\eta\wedge}"'] & {} & {} \\
{} & 0 \arrow[u] & {} & {}
\end{tikzcd}
\end{equation}

\pageoriginale

\section*{5. Los casos excepcionales del teorema de Noether}

En esta sección seguimos trabajando sobre $\mathbb C$, con cohomología
entera.

\subsubsection*{5.1.}

Sean $2n$, $d$ y $\underline a=(a_1,\ldots,a_d)$, con $d\geq1$ y
$a_i\geq2$, tales que la cohomología media de las intersecciones
completas lisas $V$ de dimensión $2n$ y multigrado $\underline a$ sea
puramente de tipo $(n,n)$. En otras palabras, o bien $n=0$, o bien
$(2n;\underline a)$ es uno de los casos $(2n;2)$, $(2n;2,2)$ y $(2;3)$
(XI 2.9).

Supongamos que $V$ se obtiene de una intersección completa
$X\subset\mathbb P^{2n+d}$, respectivamente
$X\subset\mathbb P^{2n+d+1}$, de dimensión $2n+1$ y multigrado
$(a_1,\ldots,a_{d-1})$, respectivamente $(a_1,\ldots,a_d)$, tomando
una sección por una hipersuperficie de grado $a_d$, respectivamente una
sección hiperplana, perteneciente a un pincel de Lefschetz $\ell$.

Por (4.3.1), la parte evanescente $\operatorname{Ev}(V)$ de
$H^{2n}(V)$, que coincide con su parte primitiva, está engendrada sobre
$\mathbb Z$ por los ciclos evanescentes. Es el ortogonal de la potencia
$\eta^n$ de la clase de cohomología $\eta$ de una sección hiperplana.

Como la forma de intersección sobre $H^{2n}(V)$ es unimodular y
$\sisi{\eta^{2n}}{\eta^n}$ no es $b\,y$, con $b>1$, para ninguna otra
clase de cohomología $y$ (XI 1.6(iv)), el discriminante de la forma de
intersección sobre $\operatorname{Ev}(V)$ es
\[
(\eta^n,\eta^n)=\prod_i a_i.
\]

Sea $R$ el sistema de raíces de $H^{2n}(V)$ formado por los ciclos
evanescentes y sus conjugados. Se dispone de la información siguiente:
\begin{enumerate}[label=(\alph*)]
\item $R$ es vacío o es de tipo $A$, $D$ o $E$;
\item el rango de $R$ está dado en XI 1.10;
\item su discriminante es $\prod_i a_i$. En las tablas, este
  discriminante es el orden del centro del grupo simple complejo
  simplemente conexo correspondiente. El argumento mediante el que se
  ha calculado el discriminante muestra que dicho centro, dual del
  cociente de la red de pesos por la red de raíces, es cíclico.
\end{enumerate}

\pageoriginale

En todos los casos, esta información determina completamente $R$ si
$n\ne0$. Se obtiene:

\subsubsection*{Proposición 5.2.}

\begin{enumerate}[label=(\roman*)]
\item Si $V$ es una cuádrica de dimensión par, $R$ es de tipo $A_1$.
\item Si $V$ es una superficie cúbica, $R$ es de tipo $E_6$.
\item Si $V$ tiene dimensión $2n$ y es la intersección de dos
  cuádricas, $R$ es de tipo $D_{2n+3}$.
\end{enumerate}

\subsubsection*{5.3.}

Supongamos $n\ne0$. El conjunto $R^1$ de vectores de longitud $\pm2$
en $\operatorname{Ev}(V)$ es asimismo un sistema de raíces y tiene las
mismas propiedades (a), (b) y (c) que $R$. Por tanto, $R=R^1$. Sea
$\pi_1$ el grupo fundamental del espacio de parámetros de todas las
intersecciones completas $V\subset\mathbb P^{2n+d}$ del tipo
considerado. La imagen de $\pi_1$ en
$\operatorname{GL}(\operatorname{Ev}(V))$ es el
\sisi{grupo de Weil}{grupo de Weyl} de $R$. En todos los casos es el
grupo de automorfismos de la terna
\[
\bigl(H^{2n}(V),\eta^n,\text{forma de intersección}\bigr).
\]

\subsubsection*{5.4.}

Reemplazando $R$ por $R^1$ como en 5.3, la afirmación
\sisi{4.2(ii)}{5.2(ii)} es bien conocida [1], y para valores pequeños de
$n$, \sisi{4.2(iii)}{5.2(iii)} era conocida, al menos por Manin
($n=0$ es trivial; una superficie que es intersección de dos cuádricas
es una superficie de del Pezzo). M. Reid [4] realizó un estudio
detallado de la intersección de dos cuádricas. Deduce 5.2(iii) a partir
del estudio de la geometría de los subespacios lineales
$n$-dimensionales de una intersección de dos cuádricas
$V\subset\mathbb P^{2n+2}$.

\section*{Bibliografía}

\begin{enumerate}[label={[\arabic*]}]
\item E. Cartan, ``Quelques remarques sur les 28 bitangentes d'une
  quartique plane et les 27 droites d'une surface cubique'',
  \emph{Bull. Sci. Math.} \textbf{70} (1946), 42--45;
  \emph{Oeuvres}, parte I, vol. 2, p.~1353.
\item P. Deligne, ``Théorie de Hodge II'', \emph{Publ. Math. IHES}
  \textbf{40} (1971), 5--58.

\pageoriginale

\item S. Lefschetz, \emph{L'analysis situs et la géométrie algébrique},
  Gauthier--Villars, 1924, reeditado en \emph{Selected Papers},
  Chelsea.
\item M. Reid, tesis doctoral, Cambridge, 1972.
\item J.-P. Serre, ``Faisceaux algébriques cohérents'',
  \emph{Ann. of Math.} \textbf{61} (1955), 197--278, citado como FAC.
\end{enumerate}
